\section{Marketing and Sales}

\subsection{Market Potential}

\subsubsection{Target Market}

The primary target market for this system is local government bodies in Sri Lanka, specifically municipal councils, urban councils, and pradeshiya sabhas that are responsible for the maintenance and operation of streetlights. There are hundreds of such local government bodies across Sri Lanka, and together they manage tens of thousands of streetlights. This represents a very large potential market.

\subsubsection{Market Need}

The problem validation survey confirmed that inefficient streetlight management and delayed fault detection are widespread problems. All of the organisations that manage streetlights would potentially benefit from this system.

\subsection{Financial Analysis}

\subsubsection{Cost Breakdown}

Table~\ref{tab:cost} shows the approximate cost of the components used to build one unit of the system.

\begin{table}[H]
    \centering
    \caption{Breakdown of Component Costs}
    \label{tab:cost}
    \begin{tabular}{lr}
        \toprule
        \textbf{Component} & \textbf{Approximate Price (LKR)} \\
        \midrule
        ESP32-WROOM-32E module & 1,200 \\
        HLK-20M05 power module & 850 \\
        AMS1117-3.3 voltage regulator & 50 \\
        ACS712 current sensor module & 350 \\
        LDR sensor & 30 \\
        SRD-05VDC-SL-C relay & 120 \\
        2N2222 transistor & 20 \\
        1N4007 diode & 10 \\
        EVB GSM 800L module & 1,500 \\
        Resistors and capacitors & 150 \\
        KF301-2P screw terminals & 80 \\
        PCB fabrication & 1,800 \\
        Miscellaneous (wires, SIM, etc.) & 500 \\
        \midrule
        \textbf{Total} & \textbf{6,660} \\
        \bottomrule
    \end{tabular}
\end{table}

\subsubsection{Pricing Strategy}

Based on the component cost, the system can be priced between 8,000 and 10,000 LKR per unit. This price point covers the material cost, assembly labour, and a small profit margin. Bulk orders from councils could bring the per-unit price down further.

\subsubsection{Cost Justification}

The energy savings alone from eliminating daytime operation can justify the cost of the system. A single streetlight left on unnecessarily for 8 hours per day consumes a significant amount of energy over a month. Multiplied across hundreds of streetlights, the total saving can pay back the investment in the system within a few months.

\subsection{Return on Investment (ROI)}

\subsubsection{Economic Impact}

Municipal councils currently spend money on electricity wasted by lights that run during the day, and on maintenance visits to discover and repair faults that could be detected automatically. By reducing energy waste and enabling faster fault response, the system can produce significant cost savings for the council.

\subsubsection{Payback Period}

Given the energy and maintenance savings, the payback period for the system is estimated to be less than one year for most councils. After that, the system continues to save money for as long as it remains in operation.

\subsection{Cost-Benefit Analysis}

The system provides a favourable cost-benefit ratio. The initial investment is low compared to the long-term savings from reduced energy consumption and maintenance costs. The low-cost components used in the design also mean that repairs and replacements are inexpensive. The durability of the PCB-based design ensures that the system remains operational for many years with minimal maintenance.
